% Source: https://artofproblemsolving.com/community/c3h2029806_happy_pi_day_maybe_a_competition
% Stable identity: aops-topic-2029806-q12
There are \(\binom82=28\) pairs of vertices. A pair is joined by an edge (\(L^2=64\)), a face diagonal (\(L^2=128\)) or a space diagonal (\(L^2=192\)).

\begin{asy}
usepackage("amsmath");
size(3.5cm);
defaultpen(fontsize(10pt));
pair ex=(1,0), ey=(0.45,0.35), ez=(0,1);
pair P(real x,real y,real z){return x*ex+y*ey+z*ez;}
draw(P(0,0,0)--P(1,0,0)--P(1,0,1)--P(0,0,1)--cycle,linewidth(0.9));
draw(P(1,0,0)--P(1,1,0)--P(1,1,1)--P(1,0,1)); draw(P(0,0,1)--P(0,1,1)--P(1,1,1));
draw(P(0,0,0)--P(0,1,0)--P(1,1,0),dashed); draw(P(0,1,0)--P(0,1,1),dashed);
draw(P(0,0,0)--P(1,0,0),blue+linewidth(1.2)); draw(P(0,0,0)--P(1,0,1),deepgreen+linewidth(1)); draw(P(0,0,0)--P(1,1,1),red+linewidth(1));
label("$8$",P(0.5,0,0),S,blue); label("$8\sqrt2$",P(0.3,0,0.72),deepgreen); label("$8\sqrt3$",P(1,1,0.6),E,red);
label("\mbox{}",(min(currentpicture,user=true)+max(currentpicture,user=true))/2,black); // keeps the text after the figure black when LaTeX typesets the labels inline
\end{asy}

There are \(12\) edges, \(6\cdot2=12\) face diagonals and \(4\) space diagonals (\(12+12+4=28\)). The condition \(L^2\ge100\) excludes exactly the edges, so the probability is
\[
\dfrac{12+4}{28}=\boxed{\dfrac{4}{7}}.
\]
